\section{Pluggable oracles and algorithm protocols}
\label{sec:oracles}

This section presents the mechanism that addresses R1, R2, R6, and R8: the
oracle lifecycle (\secref{sec:oracle-lifecycle}), QRAM as a first-class
resource (\secref{sec:qram}), and the protocol/contract architecture
(\secref{sec:protocols}).  \figref{fig:oracle-lifecycle} summarizes the
lifecycle; \lstref{lst:bv} shows it end-to-end on a small example.

\subsection{The oracle lifecycle: declare, carry, inspect, bind}
\label{sec:oracle-lifecycle}

\paragraph{Declare.}
An abstract oracle is declared by naming it, giving it a typed register
signature, and choosing a \emph{paradigm}: one of nine named contracts
(Table~\ref{tab:paradigms}) that fix the mathematical interface the
implementation must honor.  The declaration optionally carries resources
(typed QRAM slots), capability flags (\texttt{supports\_adjoint},
\texttt{supports\_controlled}), and attributes.  Internally the declaration
is a \texttt{Module} whose body is absent---an \emph{open module}.  Openness
is a first-class IR state, not an error state.

\begin{table}[t]
  \centering
  \caption{The nine oracle paradigms.  Each paradigm fixes the mathematical
  contract of the slot; concrete implementations are interchangeable as long
  as they honor the contract, the register-type tuple, the scale factor
  $\alpha$, and the declared capabilities.}
  \label{tab:paradigms}
  \small
  \setlength{\tabcolsep}{4pt}\begin{tabular}{@{}lp{10.4cm}@{}}
    \toprule
    \tabhead{Paradigm} & \tabhead{Contract} \\
    \midrule
    \texttt{unitary} & a unitary on the declared registers \\
    \texttt{block\_encoding} & $(\alpha, a, 0)$-block-encoding of a target operator: a unitary whose top-left block is $A/\alpha$; $\alpha$ is a slot attribute \\
    \texttt{database\_xor} & $|a\rangle|d\rangle \mapsto |a\rangle|d \oplus M[a]\rangle$ for a memory $M$ \\
    \texttt{state\_prep\_isometry} & prepares a target state (or applies an isometry) from a fixed zero input, with a clean-workspace convention \\
    \texttt{sparse\_location\_inplace} & in-place oracle for the \emph{location} structure of a sparse matrix access \\
    \texttt{sparse\_entry\_xor} & XOR oracle returning matrix \emph{entries} of a sparse access model \\
    \texttt{phase\_oracle} & applies a phase function of the input register \\
    \texttt{reversible\_function} & a reversible classical function on bit registers \\
    \texttt{algorithm\_stage} & an opaque algorithm fragment with declared capabilities (used to stage parts of solvers) \\
    \bottomrule
  \end{tabular}
\end{table}

\paragraph{Carry.}
Because an open module is a valid module, an algorithm that calls it is a
valid program: validation checks types, alias-freedom, capability demands,
and call-graph acyclicity exactly as for closed programs; serialization
round-trips the open state; structural tests run unchanged.  An algorithm
published against an open oracle slot is complete \emph{as an algorithm}
(R1), and the IR records precisely which slots remain open and why.

\paragraph{Inspect.}
Gap analysis (\texttt{unresolved}) walks the call graph breadth-first from
the entry and reports every reachable open slot with its name, paradigm,
shortest call path, register signature, and attributes.  Only reachable
slots count: an algorithm library that declares several candidate inputs but
uses only some of them is not penalized.  This makes open programs
self-describing artifacts for review, cost analysis, and tooling.

\paragraph{Bind.}
The linking pass \texttt{bind(program, bindings)} closes slots: each slot's
body is rewritten to a single call of the provided implementation, subject to
compatibility checks---exact register-type-tuple equality, paradigm match,
block-encoding factor $\alpha$ match, and possession of any capabilities the
slot advertised.  Binding is \emph{partial} (unlisted slots stay open) and
\emph{monotone} (an incompatible binding is rejected without touching the
program).  When the implementation captures QRAM resources, the linker
hoists them to entry-level resources, renames captures to collision-free
names along the call graph, and threads the resource bindings through---a
classical linker discipline, applied to a quantum IR.  Export and execution
require closure: \texttt{validate(require\_closed=True)} lists any slot that
would prevent it.  The same open program can be bound repeatedly against
different implementations (R2): a gate network for a small instance, a QRAM
table for a large one, an application circuit for production.

\begin{figure}[t]
  \centering
  \begin{tikzpicture}[
      font=\small,
      state/.style={draw=black!60, rounded corners=3pt, fill=black!3, align=center, inner sep=4pt},
      open/.style={state, fill=black!10},
      act/.style={font=\footnotesize\ttfamily, align=center},
      arr/.style={-{Stealth[length=2.2mm]}, black!65},
      lbl/.style={font=\footnotesize\itshape, black!55},
    ]
    \node[open] (decl) at (0,0)
      {\textbf{declaration}\\\footnotesize name, typed signature,\\\footnotesize paradigm, capabilities,\\\footnotesize $\alpha$, resources};
    \node[open, right=1.05cm of decl] (carry)
      {\textbf{open module}\\\footnotesize \texttt{body} $=\varnothing$; valid,\\\footnotesize serializable;\\\footnotesize algorithm is complete};
    \node[state, right=1.05cm of carry] (closed)
      {\textbf{bound module}\\\footnotesize body $=$ call of\\\footnotesize implementation; closure\\\footnotesize for export / execution};
    \node[state, below=0.75cm of carry] (inspect)
      {\textbf{gap report}\\\footnotesize open slots: paradigm,\\\footnotesize path, signature, attributes};
    \node[act, below=0.9cm of decl] (a1) {abstract\_*\\(declare)};
    \node[act, below=0.9cm of closed] (a2) {bind\\(link)};

    \draw[arr] (decl) -- node[lbl, above, midway] {carried} (carry);
    \draw[arr] (carry) -- node[lbl, above, midway] {attach} (closed);
    \draw[arr] (carry) -- (inspect);
    \draw[arr] (inspect.west) to[bend right=22] (decl.south east);
    \node[lbl] at (1.95, -3.2) {still open};
    \draw[arr] (closed.north) to[bend right=35] node[lbl, above, midway] {re-bind another realization} (carry.north);
    \draw[arr, dashed] (a1) -- (decl);
    \draw[arr, dashed] (a2) -- (closed);
  \end{tikzpicture}
  \caption{The oracle lifecycle.  Declarations create open modules that are
  carried inside otherwise-complete programs; gap analysis reports what
  remains open; the binding pass attaches implementations (partially and
  repeatedly, with compatibility checks and QRAM resource hoisting) until
  the program is closed for export or execution.}
  \label{fig:oracle-lifecycle}
\end{figure}

\begin{lstlisting}[style=pyqpython, caption={The oracle lifecycle on Bernstein--Vazirani (verbatim from the doctested tutorial).  The algorithm is written and \emph{simulable} only after binding; the same open program later binds a QRAM realization instead of the gate network.}, label={lst:bv}]
# 1. Declare the input, write the algorithm against it.
from pyqecclang.algorithms.oracles import abstract_database
from pyqecclang.algorithms.oracle_algorithms import bernstein_vazirani
from pyqecclang import unresolved

given = abstract_database("BooleanFunction", 3, 1)   # XOR database slot
opened = bernstein_vazirani(given).program()
assert [item.name for item in unresolved(opened)] == ["BooleanFunction"]

# 2. Bind a gate implementation; the program closes and runs.
from pyqecclang.algorithms.oracle_algorithms import affine_boolean_oracle
from pyqecclang import bind, simulate

implementation = affine_boolean_oracle(3, secret=5, bias=1)
closed = bind(opened, {"BooleanFunction": implementation.operation})
assert not unresolved(closed)
state = simulate(closed)
probability = sum(abs(a)**2 for k, a in state.amplitudes.items() if k[0] == 5)
assert abs(probability - 1) < 1e-12

# 3. The same open slot can instead bind a QRAM realization:
#    bind(opened, {"BooleanFunction":
#                  Binding(qram_database(3, 1).operation, {"table": "truth"})})
#    with the runtime data table supplied separately.
\end{lstlisting}

\begin{figure}[t]
  \centering
  \begin{quantikz}[remember picture, every matrix/.append style={name=pqBvAlg}]
  \lstick{$\mathtt{input}_{0}$} & \qw & \qw & \gate{H} & \gate[wires=4]{\mathtt{BooleanFunction}} & \gate{H} \\
  \lstick{$\mathtt{input}_{1}$} & \qw & \qw & \gate{H} & \qw & \gate{H} \\
  \lstick{$\mathtt{input}_{2}$} & \qw & \qw & \gate{H} & \qw & \gate{H} \\
  \lstick{$\mathtt{answer}_{0}$} & \gate{X} & \gate{H} & \qw & \qw & \qw \\
  \end{quantikz}\hfill
  \begin{quantikz}[remember picture, every matrix/.append style={name=pqBvImpl}]
  \lstick{$\mathtt{address}_{0}$} & \qw & \ctrl{3} & \qw \\
  \lstick{$\mathtt{address}_{1}$} & \qw & \qw & \qw \\
  \lstick{$\mathtt{address}_{2}$} & \qw & \qw & \ctrl{1} \\
  \lstick{$\mathtt{data}_{0}$} & \gate{X} & \targ{} & \targ{} \\
  \end{quantikz}
  \caption{The bound Bernstein--Vazirani program of \lstref{lst:bv}, drawn by
  the quantikz export.  Left: the algorithm as generated---three input wires,
  one answer wire, and the oracle as a single box named by its slot
  (\texttt{BooleanFunction}); this is the whole program, with no oracle
  internals materialized.  Right: the realization that \texttt{bind} attached
  for this run (the affine oracle with secret $5$, bias $1$): an $X$ on the
  data wire and two XOR controls from address bits $0$ and $2$.  The two
  circuits are modules of one program; the algorithm circuit is unchanged
  whichever realization closes the slot.}
  \label{fig:bv-circuit}
\end{figure}

\subsection{QRAM as a first-class resource}
\label{sec:qram}

Data access is where scientific-computing programs differ most sharply from
textbook circuits, and where the difference between an addressable memory
and an amplitude illusion matters most.  \pqlang{} makes QRAM a
\emph{typed resource} of a module: \texttt{QRAM(address\_width, data\_width)}
declared alongside registers, consumed by \texttt{Load} instructions with
precise XOR semantics~\cite{giovannetti2008qram,dalzell2025qram}, self-inverse and therefore reversible.  The data
table itself is \emph{runtime state}, not program text: it is bound at
execution (or recorded beside an export), can be patched locally between
runs, and never inflates the serialized program (R5).  Concrete
data-loading families in the library include select-swap QROM with an
explicit partition cost model, alias-sampling PREPAREs, banked multi-word
QRAM, and gate tables---all interchangeable realizations of the same oracle
slots (R2).

One subtlety deserves emphasis because it is a recurring correctness trap in
hand-written code: an XOR database of \emph{angles} is not an amplitude
oracle.  The library's diagonal block-encoding adapter, for instance,
realizes $D_{jj}=\alpha\cos(s\,d_j/2)$ from an angle table $d_j$; feeding a
fixed-point table of coefficients $c_j$ and declaring the result
$\operatorname{diag}(c_j)$ would be wrong, and the input model makes the
distinction explicit in the slot's semantics rather than in a comment.

\begin{figure}[t]
  \centering
  \begin{quantikz}[remember picture, every matrix/.append style={name=pqLookup2}]
  \lstick{$\mathtt{address}_{0}$} & \gate[wires=5]{\mathrm{QRAM}_{\mathtt{table}}} \\
  \lstick{$\mathtt{address}_{1}$} & \qw \\
  \lstick{$\mathtt{data}_{0}$} & \qw \\
  \lstick{$\mathtt{data}_{1}$} & \qw \\
  \lstick{$\mathtt{data}_{2}$} & \qw \\
  \end{quantikz}\hfill
  \begin{quantikz}[remember picture, every matrix/.append style={name=pqDbe}]
  \lstick{$\mathtt{target}_{0}$} & \gate[wires=5]{\mathtt{DB}} & \qw & \qw & \qw & \gate[wires=5]{\mathtt{DB}^{\dagger}} \\
  \lstick{$\mathtt{target}_{1}$} & \qw & \qw & \qw & \qw & \qw \\
  \lstick{$\mathtt{signal}_{0}$} & \qw & \ctrl{3} & \qw & \qw & \qw \\
  \lstick{$\mathtt{signal}_{1}$} & \qw & \qw & \ctrl{2} & \qw & \qw \\
  \lstick{$\mathtt{signal}_{2}$} & \qw & \qw & \qw & \ctrl{1} & \qw \\
  \lstick{$\mathtt{signal}_{3}$} & \qw & \gate{R_y(\frac{\pi}{4})} & \gate{R_y(\frac{\pi}{2})} & \gate{R_y(\pi)} & \qw \\
  \end{quantikz}
  \caption{Two QRAM idioms, drawn by the quantikz export.  Left: a bare
  \texttt{Load}---one instruction on the address and data registers,
  $|a\rangle|d\rangle \mapsto |a\rangle|d \oplus M[a]\rangle$, with the table
  itself a runtime resource.  Right: the diagonal block-encoding of an angle
  table ($D_{jj}=\alpha\cos(s\,d_j/2)$, here $s=\pi/4$, $\alpha=1$): query
  the angle database, apply one controlled $R_y$ per angle bit onto the flag
  wire, unquery.  The same ladder realizes the coefficient angle databases
  of the input-model matrix (\tabref{tab:input-models}); the database box is
  again a slot, bindable to a gate truth table or a QRAM resource.}
  \label{fig:qram-circuit}
\end{figure}

\subsection{Algorithm protocols and contracts}
\label{sec:protocols}

R8---new algorithms without core changes---is met by pushing all
\emph{algorithmic} typing out of the language core and into the library,
using Python's structural protocols.  The mechanism has three layers:

\paragraph{Structural interfaces.}
The library defines small \texttt{typing.Protocol} interfaces
(\texttt{runtime\_checkable}) that any object may satisfy structurally:
\texttt{UnitaryProtocol}, \texttt{Block\-Encoding\-Protocol},
\texttt{State\-Preparation\-Protocol}, \texttt{XorDatabase\-Protocol},
and \texttt{StateOracleProtocol}, among others.  The core \texttt{Operation} object
implements all applicable adapters, so a finished builder module is
simultaneously a unitary, a state preparation, and a (trivial)
block-encoding; algebra such as \texttt{lcu} consumes any objects satisfying
the right protocol, whatever their provenance.  User applications may define
their own protocols---they are ordinary Python classes with no registration
step---and check them with \texttt{requires(value, Interface)}.

\paragraph{Algorithm contracts.}
Each solver family declares a \emph{contract}: which input protocols it
requires (with widths, capabilities, and promise objects), and which output
protocols it provides.  \texttt{solver.check(problem)} produces a structured
report---issue codes such as \texttt{INPUT\_PROTOCOL},
\texttt{INPUT\_ADAPTER}, \texttt{INPUT\_WIDTH}, \texttt{INPUT\_CAPABILITY},
\texttt{INPUT\_PROMISE}, \texttt{OUTPUT\_LAYOUT}, with paths, expected and
actual values, and human-readable messages---and \texttt{.require()} turns
reports into exceptions.  Reports serialize to JSON: contracts are
machine-readable documentation, not a second executable IR.

\paragraph{Generator protocols.}
Solver families are exposed as small callables---for the ODE stack, the
three-argument \texttt{QODEProtocol} mapping (generator, initial state,
time) to a state oracle---so that \emph{solvers are plug-compatible
positions} in larger programs: a QLSS instance, an ODE solver family, or a
Hamiltonian-simulation kernel can be injected where a linear solver is
expected (\secref{sec:ode}).

\begin{lstlisting}[style=pyqpython, caption={Protocols in action (adapted verbatim from \texttt{examples/algorithm\_contracts.py}).  Top: an application-defined protocol \texttt{HasDiagonal}, satisfied structurally by an ordinary class whose block-encoding slot is an open declaration.  Bottom: a Costa-QLSS instance consumes the problem object, its contract check succeeds, the solver runs, and the resulting open program is closed by binding an implementation to the \texttt{GivenA} slot.}, label={lst:contracts}]
@runtime_checkable
class HasDiagonal(Protocol):
    """An application-local protocol; the language need not know it."""
    def diagonal_values(self) -> tuple[float, ...]: ...

class GivenMatrix:
    def __init__(self):
        self.encoding = abstract_block_encoding("GivenA", 1, 0, 1.0)

    def block_encoding(self):
        return self.encoding          # open oracle slot

    def diagonal_values(self):
        return (1.0, 1.0)

a = GivenMatrix()
assert requires(a, HasDiagonal).diagonal_values() == (1, 1)
qlss = make_costa_qlss(CostaConfig(steps=1))
problem = LinearSystem(block=BlockSystem(a, gate, SpectralPromise(1, 1)),
                       rhs_norm=1)
report = qlss.check(problem)
report.require()                       # contract check passes
result = qlss(problem)                 # -> StateOracleProtocol
closed = bind(result.operation.program(),
              {"GivenA": identity(1).operation})

limited = abstract_state_prep("ForwardOnly", 1, reversible=False)
rejected = qlss.check(LinearSystem(
    block=BlockSystem(a, limited, SpectralPromise(1, 1))))
assert not rejected.ok                 # capability violation, reported
\end{lstlisting}

\paragraph{Promise provenance.}
Mathematical promises are objects with \emph{evidence}: a spectral-promise
object (\texttt{SpectralPromise}) carries norm and smallest-singular-value
bounds together with a provenance string recording why the bound is
believed (R6).
A promise without evidence is still accepted---the language does not verify
mathematics---but the declaration travels with the program and appears in
contract reports, so certification work later knows exactly what was
assumed.  The same discipline extends to solver internals: generated modules
carry attributes recording what a kernel omitted (finite quadrature plan,
omitted contour poles, Carleman truncation), which
\secref{sec:implementation} turns into an explicit per-method verification
matrix.

\subsection{Readout as a host-side plan}
\label{sec:readout}

Because RIR is purely unitary (R7), measurement never blocks composition:
solver outputs are state oracles that further algorithms consume coherently.
When a deployment does need classical data, the host attaches
\emph{readout actions}---measure, reset, post-selection, and the statistics
to compute---as a plan appended after export, referencing registers by name.
Estimation primitives in the library (phase estimation, amplitude estimation
with its decoder, Hadamard tests, swap tests) are themselves unitary modules
whose measurement plan is standard; \secref{sec:primitives} describes them.
